\section{Compare reproducible listening streams}\label{pcm-listening}
The listening fixture uses a 64-sample integer triangle repeated at 16,000 frames
per second, giving a 250 Hz fundamental. For phase position $u$ from zero through
63 and amplitude $A$, its component values are
\[
 t(u)=\begin{cases}
 Au/16 & u<16,\\
 A(32-u)/16 & 16\leq u<48,\\
 A(u-64)/16 & 48\leq u<64.
 \end{cases}
\]
The selected amplitudes are multiples of sixteen, so these values are exact integers.
Each scene lasts 32,000 mono frames, or two seconds. The quiet scene uses amplitude
6000 in each input. The clipping scene uses 24000 in each, aligned at the same phase.
A third scene shifts the second 24000-amplitude stream by eight sample positions.

The quiet scene has no clipped samples and its saturated and wrapping output bytes
are identical. The two louder scenes each have 11,000 clipped samples under the
specified inputs, although their exact sums and output shapes differ. A clipping
count describes how many samples exceeded a boundary; it does not describe the
whole waveform or the magnitude by which those samples exceeded it.

A provided renderer feeds successive groups to the simulator, checks every output
sample against an independent integer rule, and collects the samples actually stored
by the mixing and wrapping routines. It writes ordinary sixteen-bit mono WAV files
with their sample rate and channel metadata. Both tested mixer implementations
produce identical files. There is no synthesized sound-device timing in the kernel.
A later file-format project will examine the WAV container; this project works on
the decoded PCM sample slices.

\begin{figure}[htbp]
\centering
\begin{tikzpicture}[alt={One sixty-four-sample period of the clipping listening scene compares exact sums, saturated output and wrapped output. Both inputs peak at 24000 at index sixteen and reach minus 24000 at forty-eight. Exact sums peak at 48000 and minus 48000. Saturation flattens those peaks at 32767 and minus 32768; wrapping instead gives minus 17536 and positive 17536 there. Clipping begins at index eleven, where exact 33000 wraps to minus 32536. Eleven samples in each half-period clip, for twenty-two per period.}]
\begin{axis}[width=0.88\linewidth,height=65mm,xmin=0,xmax=63,ymin=-52000,ymax=52000,xlabel={Sample index in one period},ylabel={Sample value},xtick={0,16,32,48,63},ytick={-48000,-32768,0,32767,48000},scaled y ticks=false,tick label style={font=\small},legend style={at={(0.5,-0.3)},anchor=north,legend columns=3,font=\small},grid=major]
\addplot[black,dotted] coordinates {(0,0) (1,3000) (2,6000) (3,9000) (4,12000) (5,15000) (6,18000) (7,21000) (8,24000) (9,27000) (10,30000) (11,33000) (12,36000) (13,39000) (14,42000) (15,45000) (16,48000) (17,45000) (18,42000) (19,39000) (20,36000) (21,33000) (22,30000) (23,27000) (24,24000) (25,21000) (26,18000) (27,15000) (28,12000) (29,9000) (30,6000) (31,3000) (32,0) (33,-3000) (34,-6000) (35,-9000) (36,-12000) (37,-15000) (38,-18000) (39,-21000) (40,-24000) (41,-27000) (42,-30000) (43,-33000) (44,-36000) (45,-39000) (46,-42000) (47,-45000) (48,-48000) (49,-45000) (50,-42000) (51,-39000) (52,-36000) (53,-33000) (54,-30000) (55,-27000) (56,-24000) (57,-21000) (58,-18000) (59,-15000) (60,-12000) (61,-9000) (62,-6000) (63,-3000)};
\addlegendentry{Exact sum}
\addplot[MARSAccent,thick] coordinates {(0,0) (1,3000) (2,6000) (3,9000) (4,12000) (5,15000) (6,18000) (7,21000) (8,24000) (9,27000) (10,30000) (11,32767) (12,32767) (13,32767) (14,32767) (15,32767) (16,32767) (17,32767) (18,32767) (19,32767) (20,32767) (21,32767) (22,30000) (23,27000) (24,24000) (25,21000) (26,18000) (27,15000) (28,12000) (29,9000) (30,6000) (31,3000) (32,0) (33,-3000) (34,-6000) (35,-9000) (36,-12000) (37,-15000) (38,-18000) (39,-21000) (40,-24000) (41,-27000) (42,-30000) (43,-32768) (44,-32768) (45,-32768) (46,-32768) (47,-32768) (48,-32768) (49,-32768) (50,-32768) (51,-32768) (52,-32768) (53,-32768) (54,-30000) (55,-27000) (56,-24000) (57,-21000) (58,-18000) (59,-15000) (60,-12000) (61,-9000) (62,-6000) (63,-3000)};
\addlegendentry{Saturated}
\addplot[black,dashed] coordinates {(0,0) (1,3000) (2,6000) (3,9000) (4,12000) (5,15000) (6,18000) (7,21000) (8,24000) (9,27000) (10,30000) (11,-32536) (12,-29536) (13,-26536) (14,-23536) (15,-20536) (16,-17536) (17,-20536) (18,-23536) (19,-26536) (20,-29536) (21,-32536) (22,30000) (23,27000) (24,24000) (25,21000) (26,18000) (27,15000) (28,12000) (29,9000) (30,6000) (31,3000) (32,0) (33,-3000) (34,-6000) (35,-9000) (36,-12000) (37,-15000) (38,-18000) (39,-21000) (40,-24000) (41,-27000) (42,-30000) (43,32536) (44,29536) (45,26536) (46,23536) (47,20536) (48,17536) (49,20536) (50,23536) (51,26536) (52,29536) (53,32536) (54,-30000) (55,-27000) (56,-24000) (57,-21000) (58,-18000) (59,-15000) (60,-12000) (61,-9000) (62,-6000) (63,-3000)};
\addlegendentry{Wrapped}
\end{axis}
\end{tikzpicture}
\caption{One period from the verified clipping listening fixture. Lines connect stored sample points for comparison; the kernel computes each sample independently.}
\end{figure}

Use the source-input, saturated and wrapping files together with the numerical
trace and waveform figure. Their different sample sequences can be inspected
without audio playback. The files retain the original kernel amplitudes and carry
hash receipts; the render process applies no normalization or hidden gain change.

The companion audio set includes
\href{https://raw.githubusercontent.com/nhubbard/mars-book/main/book/media/020-pcm-mixing/quiet-saturated.wav}{the quiet saturated control},
\href{https://raw.githubusercontent.com/nhubbard/mars-book/main/book/media/020-pcm-mixing/clipping-saturated.wav}{the clipped saturated mix}, and
\href{https://raw.githubusercontent.com/nhubbard/mars-book/main/book/media/020-pcm-mixing/clipping-wrapped.wav}{the wrapping comparison}.
The remaining source and phase-offset tracks accompany the same receipt. These
are separate mono files; compare corresponding scenes rather than treating the
file list as interleaved audio channels.
